% \section*{Заключение}
% \addcontentsline{toc}{section}{Заключение}

Точная оценка долгосрочных эффектов в A/B-тестировании остаётся одной
из ключевых проблем для платформ, выстраивающих процесс принятия
продуктовых решений на основе данных. Краткосрочные метрики позволяют
быстро итерироваться, но систематически не улавливают истинного
влияния интервенций на пожизненную ценность пользователя; прямое
наблюдение реализованной lLTV даёт статистически шумную оценку и
требует длительных периодов наблюдения, замедляющих продуктовый цикл.

В настоящей работе предложен фреймворк \textbf{Orbitals}~---
целеориентированная поведенческая сегментация, позволяющая оценивать
долгосрочный эффект A/B-экспериментов быстрее, надёжнее и
интерпретируемее существующих подходов. Центральная идея состоит в
отказе от ожидания реализации долгосрочных денежных потоков в пользу
анализа немедленного изменения траекторий пользователей между
устойчивыми поведенческими архетипами. Если относительные ценности
орбиталей остаются стабильными во времени, то долгосрочный эффект
можно оценивать уже в момент окончания эксперимента, не дожидаясь
полного наблюдения будущей LTV.

Тестирование на реальных данных крупной платформы электронной
коммерции подтвердило эффективность подхода. Орбитальный оцениватель
позволяет получать надёжные ранние сигналы о долгосрочных
последствиях продуктовых изменений без расширения экспериментального
окна. За счёт агрегирования на уровне орбиталей дисперсия оценки
снижается до $2.5$ раз, что конвертируется в повышенную
чувствительность на тех же объёмах выборки. При этом метод сохраняет
прозрачность: результат эксперимента описывается через конкретные
переходы пользователей между поведенческими состояниями, что
позволяет аналитикам не только видеть \textit{величину} эффекта, но и
понимать \textit{механизм} его возникновения.

Важным практическим свойством подхода является его встраиваемость в
существующие пайплайны экспериментирования. Краткосрочные метрики
продолжают выполнять роль быстрых операционных индикаторов, тогда как
орбитальный оцениватель работает как дополнительный слой анализа,
дающий ранний сигнал о долгосрочном эффекте. Это особенно ценно для
изменений, которые не приводят к немедленному росту выручки, но
смещают траектории пользователей в сторону более ценных состояний.

Вместе с тем предложенный метод не является универсальной заменой
полноценного долгосрочного наблюдения. Его корректность зависит от
устойчивости сегментации и от приближённого выполнения
мультипликативного допущения. Кроме того, для каждого нового продукта
требуется индивидуальный инжиниринг поведенческих признаков, а
допущение об условной независимости от будущих внешних факторов
(сезонности, промо-акций) должно систематически мониториться в
эксплуатации.

Перспективы дальнейших исследований видятся в нескольких направлениях.

\begin{description}
    \item[Динамические орбитали:] расширение фреймворка для поддержки
          изменяющихся во времени определений орбиталей,
          адаптирующихся к сдвигам пользовательского поведения и
          эволюции платформы.

    \item[Байесовское оценивание lLTV:] включение количественной
          оценки неопределённости в моделирование поведения
          пользователей для повышения робастности выводов, особенно
          в условиях ограниченных выборок и тяжёлых хвостов
          распределения выручки.

    \item[Интеграция с каузальной оценкой:] объединение фреймворка
          \textbf{Orbitals} с методами контрфактического обучения для
          дальнейшего повышения предсказательной точности и
          расширения класса допустимых интервенций.
\end{description}

Подводя итог, фреймворк \textbf{Orbitals} предоставляет простое,
масштабируемое и интерпретируемое решение для интеграции оценки
долгосрочных эффектов в стандартные процессы A/B-тестирования.
Фокусируясь на поведенчески осмысленных архетипах и их динамике,
метод позволяет видеть долгосрочные последствия продуктовых изменений
уже на ранних этапах эксперимента~--- и тем самым принимать решения,
которые являются не только быстрыми, но и стратегически обоснованными.
